\subsection{An independent rational check of the first block}
\label{cert:rational-small}

The optional program \path{verification/check_sjj_rational.py} checks
the $6\times6$ block on $J=\{1,3,4\}$ by exact rational arithmetic.
It uses the cosine rule \eqref{cert:rule} with $N=64$ and the sum over
$1\le a\le31$, and otherwise the rounded exponential, kernels and
basis densities specified above. There are no atom terms for these
basis inputs. The following error bound makes this smaller check
independent of the full Arb calculation.

\begin{lemma}\label{cert:rational-small-error}
For each basis input $c_n=1$ or $d_n=1$, $n\in J$, with all other data
zero, and each output node $m\in J$, the rational program computes
both entries of $\mathcal S_m(p)$ with absolute error less than
$10^{-12}$.
\end{lemma}

\begin{proof}
We separate quadrature error from rounding in the finite sum. Use the
same complex disks of radius $1/6$ as in
Lemma~\ref{cert:moment-error}. The basis density on any disk is bounded by
\[
 D_0=2\pi^2\cdot10\cdot\frac{13}{12}e^{4\pi/6}.
\]
This also bounds a simple-pole basis density. The bounds
$|\lambda|<5$, $|z|<6.2$ and $\operatorname{Im}z>-.081$ proved there
show that the transformed value and first-derivative integrands are
bounded respectively by
\[
 M_0=5D_0e^{4\pi\cdot.081}<24028,
 \qquad M_1=6.2\pi M_0<468015.
\]
Use the rational upper bound $e^x\le\mathcal E_n(x)$ from
\eqref{ind:positive-exponential}.  Put
$\overline\pi=p_*+10^{-36}>\pi$.  Here
$x=743\overline\pi/750$ lies in $[0,4)$, so the formula applies with $n=40$.
Using $M_0=(325/3)\pi^2e^{743\pi/750}$, the exact rational comparisons
\[
 \frac{325}{3}\overline\pi^2 \mathcal E_{40}(743\overline\pi/750)<24028,
 \qquad
 \frac{31}{5}\overline\pi\,
 \frac{325}{3}\overline\pi^2 \mathcal E_{40}(743\overline\pi/750)<468015
\]
prove the displayed bounds for $M_0,M_1$.  Each comparison is a fraction
inequality after clearing positive denominators; no rounded decimal
evaluation of an exponential is used.
Discrete cosine orthogonality now gives exactness through degree $63$.
The weights sum to one and each is at least $1/(64\cdot63)>0$ by
the same telescoping identity. Cauchy's estimate on the half-radius
real segments bounds the degree-$63$ Taylor tails by $M_i2^{-63}$.
The integral and quadrature have total mass one, after combining all
six pieces. Thus their discrepancies in the two moments are at most
\[
 \epsilon_0=2^{-62}M_0<5.211\cdot10^{-15},\qquad
 \epsilon_1=2^{-62}M_1<1.015\cdot10^{-13}.
\]
Using $Q_m>1.75$ and $|D_m|<2.12$, the quadrature error in either
entry of \eqref{cert:pair} is consequently less than
\begin{equation}\label{cert:rational-small-quadrature}
 \frac{2.12\epsilon_0+\epsilon_1}{1.75}<6.432\cdot10^{-14}.
\end{equation}

For completeness, two changes in the rational implementation require
care in reusing the rounding estimates. First, for a real argument $u$
the program evaluates $e_*(r)$ with $r=u\bmod2\in[0,2)$.
Periodicity is used for the exact target:
$e^{i\pi r}=e^{i\pi u}$. The approximation $p_*$ is not treated as
giving an exactly periodic exponential.  Since $r$ is real and
$|p_*r|<2(3.142)<1000$, Lemma~\ref{cert:rounded-exponential}, followed by
$|e^{ip_*r}-e^{i\pi r}|\le2\cdot10^{-36}$, gives an absolute error
less than $2\cdot10^{-32}$ for each such replacement. When $u$
contains an approximate quadrature node, first compare at that
approximate argument and then use the first-difference bound for the
exact exponential. This avoids any assumption of continuity of the
modulo operation.

Second, with $N=64$, the sum of the absolute cosine coefficients in
one weight is less than one. Including the final floor, the weight
error is therefore at most
\[
 \frac{2\cdot10^{-32}}N+10^{-40}
 <\frac{10^{-29}}N.
\]
Summing the stronger first expression, rather than the relaxed bound on
the right, gives the total weight error
\begin{equation}\label{cert:rational-small-weight-total}
 N\left(\frac{2\cdot10^{-32}}N+10^{-40}\right)
 =2\cdot10^{-32}+64\cdot10^{-40}<3\cdot10^{-32}.
\end{equation}
The node error is less than $2\cdot10^{-32}/12<2\cdot10^{-33}$.
The exact nodes are at distance at least $1/98304$ from the endpoints
of their own interval $[(j-1)/6,j/6]$, since
$1-\cos(\pi/128)\ge2/128^2$. Thus the approximate nodes remain in their
own intervals, in particular in $[0,1]$. Their transformed frequencies
have positive imaginary part even after the floor, because
$\operatorname{Im}z\ge26/435$ there.  The real-segment estimates in the
proof of Lemma~\ref{cert:rational-moment-error} also give $|z_*|<3$.
Thus every transformed call has its argument in the closed left
half-plane and satisfies $|p_*z_*m|<3.142(3)(4)<1000$.

Here is a direct propagation bound for these smaller inputs. The
real-segment estimates are $|\lambda|<3$, $|z|<3$, and basis-density
modulus less than $214$. For one piece put
\[
 A_n=\sum_{l=j}^6P_le^{i\pi t_ln},\qquad
 B_n=\sum_{l=j}^6t_lP_le^{i\pi t_ln}.
\]
Their moduli are at most ten. Replacing $b$ and the real exponentials
changes each sum by less than $3\cdot10^{-31}$; including the node
movement changes $B_n-tA_n$ by less than $7\cdot10^{-31}$.
The phase $e^{-i\pi tn}$ changes by less than $5\cdot10^{-32}$,
since $n\le4$. Substitution in \eqref{fourier:folded} therefore gives
a density error less than $3\cdot10^{-29}$.
The derivatives of the rational kernels, using $b,b_*>4/5$ and
$|t+ih|\ge2/5$, give errors less than $2\cdot10^{-32}$ in both
$\lambda$ and $z$, including their floors. The first-difference
formula for exponentials in the closed left half-plane, with $m\le4$,
then bounds the errors in
\[
 g_0(t)=\lambda(t)e^{i\pi mz(t)},\qquad
 g_1(t)=g_0(t)i\pi z(t)
\]
by $2\cdot10^{-30}$ and $3\cdot10^{-29}$, respectively.
The exact factors have moduli below $3$ and $36$; their approximations
have moduli below $4$ and $37$. Product differences bound each
derivative-integrand error by
\[
 3\cdot10^{-29}\cdot37+214\cdot3\cdot10^{-29}
 <8\cdot10^{-27};
\]
the value-integrand bound is smaller. By
\eqref{cert:rational-small-weight-total}, the total absolute weight error
is below $3\cdot10^{-32}$, and the approximate integrands have modulus
below $10^4$. Hence each rounded moment differs from its exact finite
quadrature moment by less than $10^{-26}$.

Finally, for residues $1,3,4$, the exact polynomial forms are
\[
\begin{gathered}
 Q_1=Q_3=\frac{4\pi^2}{3}(1-b),\qquad Q_4=\frac{\pi^2}{3},\\
 D_1=\frac{\pi(3+2b)}{18},\qquad
 D_3=-\frac{\pi(3+2b)}{18},\qquad D_4=\frac{7\pi b}{9}.
\end{gathered}
\]
The program replaces $\pi,b$ by $p_*,b_*$ in these expressions.
Using $|p_*^2-\pi^2|<6.3\cdot10^{-36}$, $\pi^2<10$, and $b,b_*<1$ gives
\[
\begin{aligned}
 |Q_m^\#-Q_m|
 &<\max\{(4/3)(6.3+10),\,6.3/3\}10^{-36}<3\cdot10^{-35},\\
 |D_m^\#-D_m|
 &<\max\{(5+44/7)/18,\,29/9\}10^{-36}<4\cdot10^{-36}.
\end{aligned}
\]
The exact rational lower bound used
after~\eqref{ind:pair-stability} therefore shows that both denominators
$Q_m,Q_m^\#$ exceed $7/4$. The moments have modulus below $10^4$, so the
additional division and $D_m$ errors are less than $4\cdot10^{-31}$.
Together with \eqref{cert:rational-small-quadrature}, these bounds
give an error less than $7\cdot10^{-14}<10^{-12}$ in either matrix
entry, as asserted.
\end{proof}

If every entry of this $6\times6$ matrix has error below
$\varepsilon=10^{-12}$, its induced column-norm error is below
$6\varepsilon$. The program accordingly enlarges the computed
defect norm by $\|W_\eta\|_1\,6\varepsilon$, then verifies that the
result is below $.003$ and that the resulting inverse bound is below
$32$. This is the same perturbation-of-identity argument as in
Section~\ref{cert:matrix-bounds}. It checks both small inverse
certificates without asserting that the full $N=256$ rational
prescription has been executed.
